\scanpage{61}
\deleted{\unsure{$G$}}

Groupe de galois d'une catégorie tannakienne \inserted{\unsure{graduée|gr. duée}} \inserted{($\mathbb Q_p$-linéaire)} $1/2$ simple $C$ munie d'un foncteur fibre à valeurs dans les \inserted{$F$-}cristaux \inserted{filtrés \unsure{gradués|gr. dués}} à isogénie près sur un anneau \deleted{local} \inserted{de} Val. discr. complet $V$ d'inégales car., à corps résiduel parfait $k$, corps des fractions $L$.

On a
\margin{\unsure{vect. Witt}}
\diagram{\[
\begin{array}{ccc}
W&\subset&V\\
\cap&&\cap\\
K&\subset&L
\end{array}
\]}

a) foncteur \inserted{fibre $\Phi$} vers les $F$-cristaux sur $k$ à isogénie près

$=$ certains couples $(E,F)$, où $E$ vectoriel de dim. finie sur $K$, $F$ endomorphisme $\sigma$-linéaire de $E$, tel qu'il existe un réseau invariant par $F$.

b) filtration de $\Phi\otimes_K L$.

\deleted{Le foncteur $\Phi$ \inserted{(en oubliant la $F$ structure)} \unsure{permet} d'identifier $C$ à la catégorie des repr. lin. sur $K$}

Le foncteur fondamental \inserted{(\unsure{convenable})} fournit un foncteur
\[
\unsure{T_p|V_p}:C\longrightarrow\operatorname{Mod}_f(\pi,\mathbb Q_p)
\]
où
\[
\pi=\operatorname{Gal}(\overline L/L)
\]
donc \uline{on a un foncteur fibre sur} $\mathbb Q_p$, correspondant à une représentation (à équiv. près) de $C$ comme
\[
C\simeq\operatorname{Rep}_{\mathbb Q_p}(G)
\]
$\boxed{G}$ étant un groupe pro-réductif sur $\mathbb Q_p$.
\mspage{60}
